\documentclass[11pt]{article}

\usepackage[a4paper,margin=2.3cm]{geometry}
\usepackage[T1]{fontenc}
\usepackage[utf8]{inputenc}
\usepackage{lmodern}
\usepackage{microtype}
\usepackage{amsmath,amssymb}
\usepackage{booktabs}
\usepackage{array}
\usepackage{graphicx}
\usepackage{siunitx}
\usepackage[authoryear,round]{natbib}
\usepackage{hyperref}
\usepackage{xcolor}
\usepackage{enumitem}
\usepackage{caption}

\hypersetup{
  colorlinks=true,
  linkcolor=black,
  citecolor=black,
  urlcolor=blue
}

\sisetup{detect-all=true}
\setlength{\parindent}{0pt}
\setlength{\parskip}{0.55em}

\newcommand{\figplaceholder}[2]{%
  \IfFileExists{#1}{%
    \includegraphics[width=#2]{#1}%
  }{%
    \fbox{\parbox[c][5.2cm][c]{0.92\linewidth}{\centering
      Figure file not yet committed:\\[0.4em]\texttt{\detokenize{#1}}}}
  }%
}

\title{\textbf{Physics-Guided Soft Sensing under Sparse Laboratory Measurements and Temporal Regime Shifts in Industrial Monoammonium Phosphate Production}}

\author{%
German A. Karnup\textsuperscript{a,b,*}\quad
Oleg A. Telminov\textsuperscript{b}\quad
Andrey N. Sychev\textsuperscript{c}\quad
Evgeny S. Gornev\textsuperscript{b}\\[0.35em]
\small
\href{https://orcid.org/0000-0001-6517-4712}{Karnup ORCID 0000-0001-6517-4712}\quad
\href{https://orcid.org/0000-0002-2358-3689}{Telminov ORCID 0000-0002-2358-3689}
}

\date{}

\begin{document}
\maketitle

\begin{center}
\small
\textsuperscript{a}Moscow Institute of Physics and Technology (National Research University), Dolgoprudny, Moscow Region, Russian Federation\\
\textsuperscript{b}Molecular Electronics Research Institute, JSC (NIIME), Zelenograd, Moscow, Russian Federation\\
\textsuperscript{c}UralAItech LLC, Russian Federation\\
\textsuperscript{*}Corresponding author: \textit{[e-mail to confirm before submission]}
\end{center}

\begin{quote}
\textbf{Anonymization note.}
The industrial data owner authorized use of the plant data for research calculations but requested de-identification of the company, production site, equipment identifiers, and internal signal tags. The manuscript therefore describes the facility only as an industrial monoammonium phosphate production line and uses generic variable names.
\end{quote}

\section*{Highlights}
\begin{itemize}[leftmargin=1.5em]
\item Virtual telemetry was validated on 260,642 minute records and 1,325 laboratory pairs.
\item Process-validity filtering excludes non-production conductivity excursions.
\item Simple Ridge and LSBoost models transfer better than the physics-guided model.
\item Independent post-repair validation confirms the temporal-transfer limitation.
\end{itemize}

\begin{abstract}
Industrial plants collect high-frequency telemetry while critical quality variables may be measured only every few hours. We study soft sensing for an industrial monoammonium phosphate process using 260,642 one-minute records and 1,325 paired laboratory observations. The grey-box model reconstructs latent acid properties from conductivity and combines them with material-balance calculations to estimate molar ratio and downstream density. Validation is entirely chronological: nine pre-repair rolling tests, a short-term regime-transfer test, a strict external post-repair test on 249 future laboratory pairs, and four post-repair rolling comparisons. A process-domain rule excludes conductivity values above 30, identified by plant technologists as washing or other non-production operation. Physics-guided structure improves interpretability but does not ensure temporal robustness. Across regime change and the post-repair period, Ridge and LSBoost remain substantially more accurate than the physics-guided model, demonstrating the need for regime-aware maintenance of industrial virtual telemetry.
\end{abstract}

\noindent\textbf{Keywords:} soft sensor; virtual sensor; hybrid modeling; process monitoring; regime shift; monoammonium phosphate

\section*{Nomenclature}

\begin{tabular}{@{}lll@{}}
\toprule
Symbol & Definition & Unit\\
\midrule
$c$ & online conductivity signal used as a proxy for acid state & process unit\\
$Q_A$ & clarified phosphoric-acid volumetric flow & m$^3$ h$^{-1}$\\
$\dot m_N$ & ammonia mass flow & kg h$^{-1}$\\
$Q_W$ & process-water volumetric flow & m$^3$ h$^{-1}$\\
$w_{P_2O_5}$ & mass fraction of $P_2O_5$ in the acid stream & fraction or \%\\
$\rho_A$ & phosphoric-acid density & kg m$^{-3}$\\
$R$ & $NH_3/H_3PO_4$ molar ratio & dimensionless\\
$\rho_M$ & downstream solution density & g cm$^{-3}$\\
$\tau$ & time lag between telemetry and laboratory reference & min\\
$W$ & aggregation-window length & min\\
\bottomrule
\end{tabular}

\section{Introduction}

Chemical-process plants routinely collect high-frequency online measurements while the variables that most directly describe product quality are measured less frequently by laboratory analysis or manual sampling. Soft sensors address this mismatch by inferring hard-to-measure quality variables from continuously available process measurements \citep{kadlec2009,souza2016}.

Industrial deployment is substantially harder than benchmark soft-sensor development. Real plant data contain missing values, sensor faults, short transients, asynchronous sampling, changing operating conditions, and uncertain reference measurements. Data preparation, temporal alignment, and model-maintenance strategy can therefore be as important as model class selection \citep{kadlec2009,offermans2024,dai2025}.

Hybrid and physics-guided models are often proposed as a way to improve interpretability and extrapolation by embedding empirical relationships inside process knowledge rather than learning an unconstrained input-output map \citep{sansana2021,bradley2022,schweidtmann2024,mousa2025}. Industrial soft sensors are also increasingly treated as a digitalization layer that converts existing sensing infrastructure into higher-level virtual measurements for monitoring and decision support \citep{hamid2022,pietrasik2024,boskabadi2025}.

Fertilizer production provides a relevant test case because composition and neutralization variables strongly affect product quality while several feed and intermediate-state properties are not continuously available. Hybrid soft sensing has previously been demonstrated for nutrient-content estimation in compound-fertilizer production \citep{fu2007}. The novelty claimed here is therefore not the generic application of a soft sensor to fertilizer production. Instead, this study examines an industrial monoammonium phosphate (MAP) process with four characteristics that are important for real deployment: one-minute telemetry versus laboratory measurements every few hours, latent feed properties, substantial data-quality problems, and a pronounced operating-regime shift during the observation interval.

The original engineering prototype estimated molar ratio and downstream density from online measurements and used these estimates in an offline decision-support layer. The scientific question addressed here is stricter: \emph{how robust is a physics-guided soft sensor under chronological hold-out testing, regime change, and an independent post-repair period, and how does it compare with simpler data-driven baselines?}

The contributions are:
\begin{enumerate}[leftmargin=1.8em]
\item a real industrial multi-rate soft-sensing case with 260,642 one-minute records and 1,325 paired laboratory observations;
\item a reproducible chronological evaluation using fixed 14-day training and 7-day test windows;
\item comparison of physics-guided, Ridge, physics-only, and LSBoost models on identical temporal splits;
\item explicit separation of process-validity failures from genuine model-transfer limitations;
\item a strict external post-repair validation showing the effect of long-term temporal shift.
\end{enumerate}

\section{Industrial process and data}

\subsection{Process description}

The investigated process is part of an industrial MAP production line. Clarified phosphoric acid is transferred through an intermediate stage to a tubular neutralization reactor. Ammonia is introduced for neutralization and process water is adjusted to influence downstream solution density. The principal reaction is represented in simplified form as
\begin{equation}
NH_3 + H_3PO_4 \rightarrow NH_4H_2PO_4.
\end{equation}

The two quality variables considered in this work are the $NH_3/H_3PO_4$ molar ratio $R$ and downstream solution density $\rho_M$. The operational targets used during prototype development were
\begin{equation}
R_{\mathrm{target}}=1.06,
\qquad
\rho_{M,\mathrm{target}}=1.21\ \mathrm{g\,cm^{-3}}.
\end{equation}

The reference values were obtained from sparse laboratory/operator measurements rather than at the one-minute frequency of the process historian. Two phosphoric-acid properties needed by the material-balance calculation, the $P_2O_5$ fraction and acid density, were not continuously measured with sufficient reliability and were therefore treated as latent variables.

\begin{figure}[htbp]
\centering
\fbox{\parbox[c][5.5cm][c]{0.92\linewidth}{\centering
\textbf{Figure 1 placeholder}\\[0.5em]
Redraw an anonymized process/soft-sensor schematic before public release. Show only generic units and variables: acid feed, intermediate vessel, reactor, ammonia and water feeds, downstream quality-sampling point, online telemetry, sparse laboratory references, latent-property reconstruction, material balance, and virtual outputs.}}
\caption{Anonymized process and soft-sensor architecture. Internal equipment numbers and plant tags must not appear.}
\label{fig:process}
\end{figure}

\subsection{Data sources}

The full study covers 1 March--29 September 2026 and preserves the maintenance outage from 6 June to 7 July.

\begin{table}[htbp]
\centering
\caption{Data sources used in the study.}
\label{tab:data}
\begin{tabular}{@{}p{3.0cm}p{4.3cm}p{4.4cm}p{2.7cm}@{}}
\toprule
Data source & Variables & Sampling / count & Role\\
\midrule
Plant telemetry & conductivity, acid flow, ammonia flow, water flow & 260,642 one-minute records; 184,387 jointly valid after preprocessing & model inputs\\
Laboratory/operator references & molar ratio and solution density & 1,325 paired observations & calibration and evaluation targets\\
Downstream production records & product-flow information & daily / lower-frequency & exploratory only\\
\bottomrule
\end{tabular}
\end{table}

The dataset comprises 139,679 pre-repair process rows and 120,961 post-repair rows. Laboratory coverage comprises 860 paired observations before repair and 465 paired observations after repair. For the frozen pre-repair configuration, 760 pre-repair and 249 post-repair laboratory pairs are aligned; for the adapted post-repair configuration, 256 post-repair pairs are aligned.

\subsection{Data quality and preprocessing}

The source data contain noisy process signals, short spikes and dropouts, missing or invalid values, incomplete time-series segments, and operating-state changes. The source workbook also contained 181 backward timestamp transitions and 181 duplicate timestamps; all channels were sorted together before filtering, and no observation was deleted by sorting.

The preprocessing workflow is:
\begin{enumerate}[leftmargin=1.8em]
\item read the native one-minute process series;
\item sort channels jointly by timestamp and preserve large acquisition gaps;
\item select the active redundant flow channel where applicable;
\item invalidate conductivity values above 30;
\item apply the authoritative conductivity filter and flow despiking;
\item apply process plausibility and range checks;
\item construct timestamps at which all required inputs are jointly valid;
\item associate each laboratory observation with an eligible process window determined by lag $\tau$ and window length $W$.
\end{enumerate}

A plant technologist confirmed that conductivity values above 30 are outside normal production operation and may correspond to vessel washing or other non-production tests. Accordingly, the process-validity rule $c>30$ was added to the preprocessing pipeline. This rule is process-based rather than model-error-based and is applied globally before conductivity smoothing, model fitting, and laboratory-window aggregation. Because the maintenance outage exceeds the continuity limit, temporal filters reset at gaps longer than 5 minutes; the gap-aware processing therefore produces two independent telemetry segments. None of these preprocessing steps uses laboratory target values.

\subsection{Operating-regime shift}

The conductivity trajectory changes abruptly on 21 May 2026. We define Regime~A as the period before 21 May and Regime~B as the period beginning 21 May. The boundary is used as an externally observed process change, not as a label learned by the prediction models. The maintenance outage from 6 June to 7 July is preserved explicitly and serves as the boundary for the independent post-repair validation.

\begin{figure}[htbp]
\centering
\figplaceholder{figures/fig_data_regimes_anonymized.png}{0.98\linewidth}
\caption{Process signals, invalid intervals, and the operating-regime boundary. The public figure must use generic signal names and contain no internal tag identifiers.}
\label{fig:dataregimes}
\end{figure}

\begin{figure}[htbp]
\centering
\figplaceholder{figures/fig_sparse_sampling_anonymized.png}{0.98\linewidth}
\caption{Sparse laboratory references on the process timeline, illustrating the mismatch between one-minute telemetry and laboratory sampling every few hours.}
\label{fig:sparse}
\end{figure}

\begin{figure}[htbp]
\centering
\figplaceholder{figures/fig_extended_coverage.png}{0.98\linewidth}
\caption{Extended process and laboratory coverage across the maintenance outage. The maintenance gap and the first post-maintenance data point are shown explicitly. For public release, the upper-panel internal signal tag must be replaced by a generic conductivity-signal label.}
\label{fig:extendedcoverage}
\end{figure}

\section{Methodology}

The overall soft-sensor architecture is summarized in Fig.~\ref{fig:soft_sensor_architecture}. The main solid-arrow pathway represents online inference. Minute-scale process telemetry first passes process-validity filtering and temporal alignment; the conductivity signal is then used to reconstruct latent acid properties, which are propagated through the physics-guided material-balance model to obtain the virtual molar-ratio and solution-density estimates. The sparse laboratory measurements are not direct runtime inputs. Instead, they are used only for calibration, chronological model selection, and validation, as indicated by the dashed training pathway. This separation between the online inference path and the calibration path is important for interpreting the model as a deployable virtual-sensing layer rather than as an interpolation of contemporaneous laboratory measurements.

\begin{figure}[htbp]
\centering
\figplaceholder{figures/soft_sensor_architecture.pdf}{0.98\linewidth}
\caption{Architecture of the physics-guided soft sensor and its calibration/validation workflow. Solid arrows denote the online inference pathway, whereas dashed arrows denote the training, calibration, and validation pathway based on sparse laboratory reference measurements.}
\label{fig:soft_sensor_architecture}
\end{figure}

\subsection{Physics-guided soft sensor}

The soft sensor is a grey-box sequence rather than a direct unconstrained regression. Conductivity is first mapped to a latent $P_2O_5$ fraction,
\begin{equation}
\widehat{w}_{P_2O_5}=f_{\theta}(c),
\end{equation}
followed by a mapping from estimated composition to acid density,
\begin{equation}
\widehat{\rho}_{A}=g_{\phi}\!\left(\widehat{w}_{P_2O_5}\right).
\end{equation}
The candidate mappings are linear, quadratic, and PCHIP spline functions.

The acid mass flow and $P_2O_5$ mass flow are
\begin{equation}
\dot m_A=\widehat{\rho}_A Q_A,
\end{equation}
\begin{equation}
\dot m_{P_2O_5}=\dot m_A\widehat{w}_{P_2O_5}.
\end{equation}

The prototype converts $P_2O_5$ mass to equivalent 100\% phosphoric-acid mass using
\begin{equation}
\dot m_{H_3PO_4,100\%}=1.38\,\dot m_{P_2O_5}.
\end{equation}

The corresponding molar flows are
\begin{equation}
\dot n_{H_3PO_4}=
\frac{\dot m_{H_3PO_4,100\%}}{M_{H_3PO_4}},
\qquad
\dot n_{NH_3}=
\frac{\dot m_N}{M_{NH_3}},
\end{equation}
with $M_{H_3PO_4}=98$ g mol$^{-1}$ and $M_{NH_3}=17$ g mol$^{-1}$. The virtual molar ratio is
\begin{equation}
\widehat R=
\frac{\dot n_{NH_3}}{\dot n_{H_3PO_4}}.
\end{equation}

For publication, the downstream density balance is written in dimensionally explicit form as
\begin{equation}
\widehat{\rho}_M=
\frac{\sum_i \dot m_i}{\sum_i Q_i},
\qquad
\dot m_i=\rho_iQ_i,
\end{equation}
where the sums contain the streams used by the implemented downstream balance. This expression uses the additive-volume approximation of the prototype. Any plant-specific correction terms should be stated separately rather than folded into the definition of density.

\subsection{Time alignment and model selection}

Candidate lags and aggregation windows are
\begin{equation}
\tau\in\{0,10,\ldots,180\}\ \mathrm{min},
\qquad
W\in\{10,20,30,60\}\ \mathrm{min}.
\end{equation}

For each validation scenario, lag, window, and physics-guided mapping type are selected using only the corresponding training data. The first 9 days of a training interval are used for inner fitting and the following 5 days for inner chronological validation. The selection score is
\begin{equation}
J=
\frac{\mathrm{RMSE}_R}{0.03}
+
\frac{\mathrm{RMSE}_{\rho}}{0.01}.
\end{equation}

The selected configuration is then frozen for the relevant outer scenario, while fitted coefficients are re-estimated from each outer training interval only.

For the final pre-repair frozen model, the selected configuration is $\tau=0$ min, $W=60$ min, with a quadratic mapping. For the adapted post-repair physics-guided model, the selected configuration is $\tau=140$ min, $W=10$ min, with PCHIP spline mapping.

\subsection{Baseline models}

All baselines use the same temporally aligned rows as the physics-guided model.

\begin{itemize}[leftmargin=1.5em]
\item \textbf{Ridge}: standardized multivariate Ridge regression using conductivity, acid flow, ammonia flow, and water flow; regularization selected on chronological inner validation.
\item \textbf{Physics-only}: acid flow, ammonia flow, and water flow with $P_2O_5$ fixed at 52 wt\% and an unfitted acid-density prior.
\item \textbf{Gradient boosting}: LSBoost with 100 cycles, learning rate 0.05, minimum leaf size 10, and maximum 20 splits.
\end{itemize}

The physics-only model is a reference calculation, not a separately calibrated mechanistic model.

\subsection{Chronological validation design}

Validation is entirely temporal. No random shuffling is used.

\textbf{Pre-repair evaluation.} The prespecified outer protocol is a fixed 14-day training interval followed immediately by a 7-day held-out interval. Within Regime~A, rolling origins advance by 7 days and yield nine complete tests (A01--A09). Regime~B does not contain the 21 days of paired laboratory coverage required for a complete within-regime 14+7-day split.

\textbf{Short-term regime-transfer evaluation.} Cross-regime transfer uses 7 May--21 May 2026 for training and 21 May--28 May 2026 for testing.

\textbf{External post-repair evaluation.} The strict external test trains on all eligible pre-repair observations and evaluates on the post-repair period 7 July--29 September 2026. This scenario is denoted \texttt{external\_transfer\_full}.

\textbf{Post-repair frozen-versus-adapted comparison.} Ten post-repair calendar windows are considered, each using 14 days for training and the next 7 days for testing. A window enters the pooled comparison only if temporal alignment leaves at least 40 training and 10 common test laboratory pairs. Four of the ten candidate windows satisfy this rule (P01--P04); the other six are excluded explicitly and do not contribute to pooled metrics.

\subsection{Evaluation metrics}

For target $y_i$ and prediction $\hat y_i$,
\begin{equation}
\mathrm{MAE}=\frac{1}{n}\sum_{i=1}^{n}|y_i-\hat y_i|,
\end{equation}
\begin{equation}
\mathrm{RMSE}=
\sqrt{\frac{1}{n}\sum_{i=1}^{n}(y_i-\hat y_i)^2},
\end{equation}
\begin{equation}
\mathrm{Bias}=
\frac{1}{n}\sum_{i=1}^{n}(\hat y_i-y_i).
\end{equation}

We also report P95 absolute error and the fractions within
\begin{equation}
|R-\widehat R|\le0.03,
\qquad
|\rho_M-\widehat{\rho}_M|\le0.01\ \mathrm{g\,cm^{-3}}.
\end{equation}

Molar-ratio predictions from the physics-guided implementation are rounded to two decimal places before validation, matching the production reporting setting.

\section{Results}

\subsection{Within-regime-A rolling validation}

Nine non-overlapping 7-day test intervals produced 499 held-out laboratory observations in total. Training-set size varied from 96 to 138 observations.

\begin{table}[htbp]
\centering
\caption{Pooled held-out performance across the nine Regime-A test intervals.}
\label{tab:pooledA}
\resizebox{\linewidth}{!}{%
\begin{tabular}{@{}lrrrrrr@{}}
\toprule
Model & RMSE $R$ & MAE $R$ & Within $\pm0.03$ & RMSE $\rho_M$ & MAE $\rho_M$ & Within $\pm0.01$\\
\midrule
Physics-guided & 0.0362 & 0.0267 & 0.581 & 0.0111 & 0.0086 & 0.667\\
Ridge & 0.0155 & 0.0114 & 0.914 & 0.0085 & 0.0063 & 0.830\\
Physics-only & 0.1055 & 0.1005 & 0.020 & 0.0841 & 0.0837 & 0.000\\
Gradient boosting & 0.0171 & 0.0130 & 0.886 & 0.0110 & 0.0086 & 0.629\\
\bottomrule
\end{tabular}}
\end{table}

The process-validity update removes the previously observed catastrophic behavior in split A08 without discarding the complete test interval. A08 retains 37 held-out laboratory observations and yields RMSE 0.0419 for molar ratio and 0.0115 g cm$^{-3}$ for density. The threshold was introduced after process review identified conductivity above 30 as non-production operation, then applied globally to the full dataset before preprocessing and model fitting.

Across Regime A, the physics-guided model now shows bounded but consistently lower predictive accuracy than Ridge. Ridge regression is the most stable model across the rolling tests for both outputs jointly. Gradient boosting gives similarly low molar-ratio errors but slightly higher density error than Ridge. The physics-only reference is systematically biased and does not approach the specified tolerances.

\begin{figure}[htbp]
\centering
\figplaceholder{figures/fig_predicted_vs_measured.png}{0.96\linewidth}
\caption{Held-out measured versus predicted values after application of the production-validity filter.}
\label{fig:scatter}
\end{figure}

\subsection{A-to-B regime transfer}

The transfer experiment trains on the final 14 days of Regime A and tests on the first 7 days of Regime B, with 84 training and 56 test laboratory observations.

\begin{table}[htbp]
\centering
\caption{Held-out A-to-B transfer performance. Density errors are in g cm$^{-3}$.}
\label{tab:transfer}
\resizebox{\linewidth}{!}{%
\begin{tabular}{@{}lrrrrrrrr@{}}
\toprule
Model & RMSE $R$ & MAE $R$ & Bias $R$ & Within $\pm0.03$ & RMSE $\rho_M$ & MAE $\rho_M$ & Bias $\rho_M$ & Within $\pm0.01$\\
\midrule
Physics-guided & 0.0966 & 0.0932 & -0.0932 & 0.018 & 0.0588 & 0.0574 & +0.0574 & 0.000\\
Ridge & 0.0170 & 0.0132 & +0.0093 & 0.875 & 0.0130 & 0.0106 & +0.0082 & 0.536\\
Physics-only & 0.0966 & 0.0912 & +0.0912 & 0.018 & 0.0877 & 0.0871 & -0.0871 & 0.000\\
Gradient boosting & 0.0174 & 0.0136 & +0.0039 & 0.911 & 0.0134 & 0.0110 & -0.0040 & 0.536\\
\bottomrule
\end{tabular}}
\end{table}

The transfer test changes the interpretation of the prototype. After the production-validity filter has removed non-production excursions, the physics-guided architecture still generalizes substantially worse than the data-driven baselines across the regime boundary. Ridge gives the lowest molar-ratio and density RMSE values, while gradient boosting is nearly identical in RMSE and gives the highest molar-ratio tolerance hit rate. The physics-guided model exhibits a strong negative molar-ratio bias and positive density bias after the regime transition.

\begin{figure}[htbp]
\centering
\figplaceholder{figures/fig_holdout_timeseries.png}{0.96\linewidth}
\caption{A-to-B transfer hold-out: physics-guided predictions and residuals.}
\label{fig:transfer_timeseries}
\end{figure}

\begin{figure}[htbp]
\centering
\figplaceholder{figures/fig_baseline_comparison.png}{0.96\linewidth}
\caption{Held-out baseline comparison. RMSE is normalized by the corresponding engineering tolerance.}
\label{fig:baseline}
\end{figure}

\subsection{No complete within-regime-B test}

The first paired laboratory record in Regime B occurs after the 21 May boundary and the final paired record is on 5 June. Consequently, there is no complete 14-day training plus 7-day held-out interval entirely within Regime B. The study reports this absence directly rather than shortening the test horizon or introducing a different protocol after observing the data.

\subsection{Independent external post-repair validation}

The most stringent evaluation uses all eligible pre-repair observations for training and tests the frozen models on the future post-repair period from 7 July to 29 September 2026. This deployment-style experiment contains 760 aligned training pairs and 249 aligned post-repair test pairs.

\begin{table}[htbp]
\centering
\caption{Strict external post-repair validation (\texttt{external\_transfer\_full}).}
\label{tab:externaltransfer}
\resizebox{\linewidth}{!}{%
\begin{tabular}{@{}lrrrrrrrr@{}}
\toprule
Model & RMSE $R$ & MAE $R$ & Bias $R$ & Within $\pm0.03$ & RMSE $\rho_M$ & MAE $\rho_M$ & Bias $\rho_M$ & Within $\pm0.01$\\
\midrule
Physics-guided frozen & 0.1277 & 0.1023 & -0.0931 & 0.112 & 0.0335 & 0.0289 & +0.0286 & 0.076\\
Ridge frozen & 0.0321 & 0.0174 & -0.0049 & 0.775 & 0.0147 & 0.0116 & +0.0104 & 0.490\\
Physics-only & 0.0962 & 0.0685 & +0.0089 & 0.281 & 0.0651 & 0.0637 & -0.0635 & 0.000\\
Gradient boosting frozen & 0.0328 & 0.0196 & -0.0098 & 0.711 & 0.0137 & 0.0102 & +0.0060 & 0.602\\
\bottomrule
\end{tabular}}
\end{table}

This result is the strongest evidence in the paper. Because neither fitting nor hyperparameter selection uses post-repair target values, the table represents a true future-period validation. The physics-guided model preserves interpretability but transfers poorly; Ridge and LSBoost remain substantially more accurate on both outputs.

\subsection{Post-repair frozen versus adapted comparison}

The post-repair rolling comparison evaluates frozen and adapted models on identical held-out laboratory rows across four valid windows (P01--P04), for a pooled total of 95 test pairs.

\begin{table}[htbp]
\centering
\caption{Post-repair rolling comparison on identical held-out samples (4 included windows, 95 pooled test pairs).}
\label{tab:postrepaircomparison}
\resizebox{\linewidth}{!}{%
\begin{tabular}{@{}lrrrrrrrr@{}}
\toprule
Model & RMSE $R$ & MAE $R$ & Bias $R$ & Within $\pm0.03$ & RMSE $\rho_M$ & MAE $\rho_M$ & Bias $\rho_M$ & Within $\pm0.01$\\
\midrule
Physics-guided frozen & 0.1386 & 0.1218 & -0.1218 & 0.021 & 0.0366 & 0.0331 & +0.0331 & 0.021\\
Ridge frozen & 0.0195 & 0.0141 & -0.0061 & 0.811 & 0.0136 & 0.0105 & +0.0102 & 0.537\\
Physics-only & 0.0753 & 0.0574 & -0.0136 & 0.274 & 0.0622 & 0.0610 & -0.0610 & 0.000\\
Gradient boosting frozen & 0.0233 & 0.0179 & -0.0124 & 0.705 & 0.0139 & 0.0097 & +0.0067 & 0.663\\
Physics-guided adapted & 0.1290 & 0.1129 & +0.0111 & 0.105 & 0.0234 & 0.0198 & -0.0040 & 0.274\\
Ridge adapted & 0.0235 & 0.0182 & -0.0039 & 0.747 & 0.0116 & 0.0087 & +0.0032 & 0.674\\
Gradient boosting adapted & 0.0260 & 0.0203 & -0.0001 & 0.663 & 0.0127 & 0.0104 & -0.0057 & 0.568\\
\bottomrule
\end{tabular}}
\end{table}

Adapting the physics-guided model on the new period improves density relative to the frozen physics-guided version but does not recover competitiveness with Ridge or LSBoost. In contrast, the data-driven models remain substantially more stable.

\begin{figure}[htbp]
\centering
\figplaceholder{figures/fig_extended_adaptation_comparison.png}{0.98\linewidth}
\caption{Post-repair held-out performance on identical rolling test samples. Bars show RMSE normalized by the engineering tolerance.}
\label{fig:postrepairrmse}
\end{figure}

\begin{figure}[htbp]
\centering
\figplaceholder{figures/fig_extended_transfer_timeseries.png}{0.98\linewidth}
\caption{Frozen pre-repair physics-guided model versus 14-day rolling adaptation on post-repair held-out dates. The adapted model reduces density bias but does not resolve the molar-ratio transfer limitation.}
\label{fig:frozenadapted}
\end{figure}

\subsection{Implications for model selection}

Across all validations, the hold-out results do not support a claim that the current physics-guided model is the most accurate virtual sensor. Instead, they distinguish two different issues. First, non-production measurements can create misleading apparent failures if process-state validity is not enforced. Second, even after invalid intervals are excluded, the physics-guided calibration transfers poorly across operating-regime change and the post-repair period.

For the present dataset, Ridge is the most stable model across both outputs. LSBoost is generally close to Ridge and sometimes slightly better on a single target, while the physics-only reference remains substantially worse than all fitted models.

\section{Discussion}

\subsection{What problem is solved?}

The study solves a practical digitalization problem: it shows how to reconstruct infrequently measured quality variables from minute-scale industrial telemetry in a MAP production process and, more importantly, how to evaluate whether such virtual telemetry remains trustworthy over time.

The contribution is therefore not only a specific soft sensor. It is a validation framework for hybrid industrial soft sensing under sparse laboratory supervision, combining process-domain filtering, chronological train/test design, explicit regime-transfer testing, and an independent post-repair future-period evaluation.

\subsection{Physics guidance does not guarantee robustness}

Physics-guided structure provides interpretability: conductivity is used to reconstruct latent acid composition and density, and these variables enter a material-balance calculation. However, the learned latent-property mappings are still empirical. When their calibration becomes unstable or the conductivity-to-composition relationship changes, the physical downstream equations propagate rather than remove the error.

The comparison with Ridge and gradient boosting is therefore not an argument against physics-guided modeling in general. It shows that the present grey-box implementation needs explicit safeguards: bounded latent-property mappings, extrapolation detection, regime-conditioned calibration, or fallback to a more stable predictor when the current input domain is outside calibration support.

\subsection{Regime shift and maintenance}

The A-to-B experiment demonstrates that the relation among conductivity, latent feed properties, and downstream quality is not stationary enough for the current calibration to be carried unchanged across the boundary. A deployment architecture should include operating-regime detection, drift monitoring, bounded mappings where physically justified, controlled recalibration, and fallback behavior selected on prospective validation.

\subsection{Why the post-repair extension strengthens the article}

The extended dataset changes the manuscript from a short-horizon industrial case study to a stronger temporal generalization study. The external post-repair experiment is especially important because it evaluates the frozen pre-repair models on a completely future period without using post-repair targets. This makes the central conclusion much more robust: the relative weakness of the physics-guided model is not an artifact of one short interval around 21 May.

\subsection{Sparse labels and validation design}

The process historian is large in minute-level telemetry count, but the number of independent target observations is much smaller. This makes random row-wise splitting inappropriate because adjacent process observations are highly correlated and the laboratory sampling process is sparse and multi-rate. The fixed 14-day/7-day protocol provides a more operationally meaningful estimate of future performance.

The post-repair rolling comparison further illustrates that sample eligibility after alignment matters. Only four of ten calendar windows satisfy the pre-specified minimum of 40 aligned training and 10 common test pairs. Explicitly reporting the excluded windows avoids hidden selection bias.

\subsection{Preprocessing and operational look-ahead}

The validation contains no outer-test target leakage: hyperparameter selection is nested chronologically and outer-test laboratory values are not used for fitting. Two preprocessing details nevertheless matter for prospective deployment.

First, the authoritative conductivity filtering and signal-repair routines include centered operations, including robust windows, Savitzky--Golay smoothing, and PCHIP repair, which can use nearby future process values. Second, laboratory alignment uses a centered averaging window around $t_{\mathrm{lab}}-\tau$; when $\tau<W/2$, the window can extend beyond the nominal laboratory timestamp.

These operations are retained because the present study validates the implemented R\&D prototype. A truly online deployment study should replace them with causal alternatives and repeat the validation.

\subsection{Decision-support scope}

The prototype also calculates 10-minute recommendations for acid, ammonia, and water flows. Those calculations remain an offline decision-support demonstration. Because recommendation quality inherits soft-sensor error and no prospective plant intervention was performed, the present article does not claim closed-loop control performance, safety benefit, or economic benefit.

The production-output calculation used during the R\&D project is likewise excluded from the primary claim because contemporaneous density measurements were unavailable for part of that calculation. The plant uses a moisture coefficient $K=1.55$; that module is retained as Supporting Information until independently validated with complete time-series inputs.

\section{Limitations}

\begin{enumerate}[leftmargin=1.8em]
\item The plant-specific process schematic and some figure labels still require final public anonymization.
\item The authoritative preprocessing contains short-horizon filtering operations; a fully causal deployment implementation should be validated separately.
\item The maintenance outage creates only four post-repair windows that satisfy the minimum-sample rule for the frozen-versus-adapted comparison.
\item Raw industrial data cannot be released publicly under the data-owner agreement.
\item The duplicate flow channel F315-2 was assessed only diagnostically and is not yet a validated failover input.
\end{enumerate}

\section{Conclusions}

A physics-guided soft-sensor prototype was evaluated for industrial MAP production using 260,642 minute-scale process records and 1,325 paired laboratory measurements. A process-domain rule excluded conductivity values above 30 as non-production operation, removing spurious failures caused by washing or test states. Chronological validation then showed that the main residual weakness of the physics-guided approach is temporal transfer rather than boundedness within the production domain.

Across nine pre-repair rolling windows, the physics-guided model remained interpretable but less accurate than Ridge. In the short-term A-to-B regime transfer test, its performance degraded substantially. The strongest evidence came from the independent post-repair validation: when trained only on pre-repair data, the physics-guided model achieved RMSE values of 0.1277 for molar ratio and 0.0335 g cm$^{-3}$ for density on 249 future laboratory pairs, whereas Ridge achieved 0.0321 and 0.0147, respectively. Rolling post-repair adaptation improved the physics-guided density estimate but did not close the gap to the data-driven baselines.

The study therefore shows that physics-guided virtual telemetry can provide interpretable estimates, but physical structure alone does not guarantee robustness to operating change. Industrial deployment should combine process-state validation, leakage-safe temporal testing, regime-aware maintenance, and explicit fallback behavior under sparse laboratory supervision.

\section*{Author contributions}

\textbf{To be confirmed before submission.} Final CRediT roles must be agreed by all authors. Relevant roles include Conceptualization, Methodology, Software, Validation, Formal analysis, Investigation, Data curation, Writing -- original draft, Writing -- review \& editing, Visualization, Project administration, and Supervision.

\section*{Funding}

This work was performed within an industrial R\&D project. The industrial data owner requested anonymization of the company, production site, equipment identifiers, and internal project identifiers. Funding metadata should be finalized before journal submission in a form consistent with that agreement.

\section*{Declaration of competing interest}

The authors declare that they have no known competing financial interests or personal relationships that could have appeared to influence the work reported in this paper.

\section*{Data availability}

The industrial data owner authorized use of the operational data for research calculations but requested de-identification of the company, site, equipment identifiers, and internal signal tags. Raw historian and laboratory data therefore cannot be made public. The manuscript reports derived aggregate metrics and the validation protocol. Subject to the same confidentiality constraints, analysis code and de-identified derived results may be released with the final article.

\section*{Declaration of generative AI and AI-assisted technologies}

Generative AI tools were used for language and structural assistance during manuscript preparation. All scientific content, equations, calculations, citations, interpretations, and conclusions are subject to author verification and approval. The final disclosure should be adapted to the policy of the selected journal at submission.

\section*{Supporting Information}

Recommended Supporting Information:
\begin{itemize}[leftmargin=1.5em]
\item exact NTO preprocessing thresholds and signal-validity rules;
\item the exact list of included and excluded post-repair windows;
\item all per-split validation metrics;
\item latent-property model coefficients for each chronological fit;
\item all per-split validation metrics;
\item validation of the production-domain conductivity exclusion and its effect on A08;
\item baseline configurations and regularization settings;
\item additional residual distributions;
\item diagnostic comparison of F315 and F315-2;
\item decision-support equations and offline examples;
\item production-output reconstruction using the plant moisture coefficient $K=1.55$;
\item data-provenance table using anonymized variable names.
\end{itemize}

\bibliographystyle{plainnat}
\bibliography{references}

\end{document}